\documentclass[10pt,a4paper]{amsart}
\usepackage[margin=19mm]{geometry}
\usepackage{amsmath,amssymb,mathtools}
\usepackage[scr=rsfs,cal=euler]{mathalfa}
\usepackage{xfrac}
\usepackage[unicode,hidelinks,bookmarks=false]{hyperref}
\newcommand{\ie}{i.e.\ }
\newcommand{\cf}{cf.\ }
\newcommand{\resp}{respectively\ }
\newcommand{\Subsection}{\subsection}
\providecommand{\nobreakdash}{\nobreak\hbox{-}}
\setlength{\emergencystretch}{2em}
\multlinegap0pt
\usepackage{tikz}
\usepackage{amssymb}
\usepackage{cancel}
\newtheorem{corollary}{Corollary}
\newtheorem{theorem}{Theorem}
\usepackage{cleveref}
\crefname{corollary}{Corollary}{Corollarys}
\crefname{theorem}{Theorem}{Theorems}
\title{\vspace*{-.4in} {Causal Geodesy: Counterfactual Estimation Along the Path Between Correlation and Causation}}
\author{Kyle Schindl, Larry Wasserman}
\date{Source: arXiv:2508.08499v2 (CC BY 4.0)}
\begin{document}
\maketitle

\noindent\emph{Source and licence.} \vspace*{-.4in} {Causal Geodesy: Counterfactual Estimation Along the Path Between Correlation and Causation}. Version arXiv:2508.08499v2 (CC BY 4.0), distributed by arXiv under the Creative Commons Attribution 4.0 International licence, http://creativecommons.org/licenses/by/4.0/, as recorded in the arXiv licence metadata for this exact version and shown on the abstract page https://arxiv.org/abs/2508.08499v2. Full licence notice: \texttt{licenses/arxiv-2508.08499-CC-BY-4.0.txt}.\par\smallskip

\noindent\textit{Editorial scope.} Counted unit: the article's theorem wasserstein\_asymptotics (source0.tex lines 570--582). The article's complete original proof of it is delivered with it (source0.tex lines 2008--2291, one \texttt{proof} environment of 284 lines, which the article places away from the statement). No mathematics is written, repaired or re-derived: the statement and the proof are the article's own lines, in the article's own notation, and no retained line is substituted or re-flowed.  Retained uncounted as the source-local context the proof needs: lines 527--537; lines 697--699; lines 707--722.  Nothing was omitted from any retained span, so the package carries no omission record.  No retained line contains a citation, so the wrapper carries no bibliography and no bibliography entry.  No retained line contains a figure environment, an image-inclusion command or drawing code, so no image is delivered and the package declares no asset dependency.  Editorial formatting only, in addition: an AMS \texttt{amsart} preamble that restates the article's own declarations and macros used by the retained lines, namely the environments \texttt{corollary}, \texttt{theorem} on separate counters, together with the standard packages those lines need.  The retained environments print in this document as Corollary 1, Theorem 1 in order of appearance, whereas the article prints the same items with the numbering of its own full document.  The complete native source of the article as distributed in the arXiv e-print for this exact version is bundled unchanged as \texttt{sources/183106/source0.tex}.
\begin{corollary}[Nonparametric Efficiency Bound] \label{var_corollary}
    Let $\sigma^2_\nu(t) = \mathbb{V}(\varphi_\nu(Z; t))$. For $t \in [0, 1)$, suppose that $0 < \pi_{\min} \leq \pi(a \! \mid \! x) \leq \pi_{\max} < \infty$ and $\sigma^2(x, a) \geq \sigma^2_{\min} > 0$ and $\mathbb{E}[Y^2 \mid X = x, A = a] \leq \sigma^2_{\max}$ for all $x \in \mathcal{X}$ and $a \in \mathcal{A}$.  Then, 
    \begin{align*}
    \sigma^2_\nu(t) = \mathbb{E}\! \left[ \left(\frac{\nu_t(A\mid X)}{\pi(A\mid X)}\right)^2 \sigma^2(X, A)  \right] + \mathbb{V}\big(\mu\big(X, (1 - t)A + ta_\star\big)\big) \asymp \frac{1}{1-t} .
    \end{align*}
    Furthermore, suppose for $P_X$-almost every $x$ that $a \mapsto \pi(a \! \mid \! x)$ and $a \mapsto \mathbb{V}(Y \mid X = x, A=a)$ are continuous at $a_\star$. Then, as $t \to 1$,
    \begin{align*}
        (1-t)\sigma^2_\nu(t) \longrightarrow  \mathbb{E}\! \left[\frac{\sigma^2_\star(X)}{\pi_\star(X)} \| \pi(\, \cdot \mid X) \|^2_{L_2(\mathcal{A})} \right],
    \end{align*}
where we say $\sigma^2_\star(X) = \mathbb{V}(Y \mid X, A = a_\star)$ and $\pi_{\star}(X) = \pi(a_\star \mid X)$.
\end{corollary}

\begin{align} \label{no_flux_assumption}
    \lim_{a\downarrow 0} \frac{\partial^j}{\partial a^j} \left\{ (a_\star-a)^k\pi(a\mid x) \right\} =  \lim_{a\uparrow 1} \frac{\partial^j}{\partial a^j} \left\{ (a_\star-a)^k\pi(a\mid x) \right\} = 0
\end{align}

\begin{theorem}[Efficient Influence Function for Wasserstein Derivatives] \label{derivative_EIF}

Let $k\geq1$ be a fixed integer and let $t \in [0, 1)$. Suppose that, for $P_X$-almost every $x$, $a\mapsto\mu(x,a)$ is $k$ times weakly differentiable and for $u \in \mathcal{A}_t$, $u \mapsto \nu_t(u\mid x)(\frac{a_\star-u}{1-t})^k$ has $k$ derivatives on $\mathcal{A}_t$. Furthermore, assume the no-flux boundary condition of \cref{no_flux_assumption} holds. Then the efficient influence function of $\psi^{(k)}_\nu(t)$ under a nonparametric model is given by
\begin{align*}
       \varphi_{\nu,k}(Z;t) =  \frac{\eta_{k, t}(X, A)}{\pi(A \mid X)}\{Y-\mu(X,A)\} + (a_\star-A)^k \frac{\partial^k}{\partial u^k}\mu(X,u) \Big|_{u=(1-t)A + ta_\star} - \psi^{(k)}_\nu(t),
\end{align*}
where 
\begin{align*}
    \eta_{k,t}(x,a) = \frac{\mathbb{I}(a\in\mathcal{A}_t)}{(1-t)^{k+1}} \sum_{j=0}^{k} \binom{k}{j} \frac{k!}{j!} \left( \frac{a-a_\star}{1-t}\right)^j \pi^{(j)}\left(\frac{a-ta_\star}{1-t}\mid x\right).
\end{align*}
Furthermore, the nonparametric efficiency bound is given by
\begin{align*}
    \mathbb{V}(\varphi_{\nu,k}(Z;t)) = \mathbb{E} \! \left[ \int_{\mathcal{A}_t} \frac{\eta_{k,t}(X, a)^2}{\pi(a\mid X)} \sigma^2(X,a)da \right] + \mathbb{V} \!\left[(a_\star-A)^k \left. \frac{\partial^k}{\partial u^k}\mu(X,u) \right|_{u=(1-t)A + t a_\star} \right].
\end{align*}
    
\end{theorem}

\begin{theorem} \label{wasserstein_asymptotics}
    Suppose that $|Y| \leq C$ with probability one, the conditions of \cref{var_corollary} hold, and that $0 < \pi_{\min} \leq \widehat{\pi}(a \! \mid \! x) \leq \pi_{\max}$ for all $a \in \mathcal{A}$ and $x \in \mathcal{X}$. Then, as $t_n \to 1$, if
    \begin{enumerate}
        \item[$(i)$] $\sqrt{n(1-t_n)} \to \infty$.
        \item[$(ii)$]  $ \|\widehat{\pi} - \pi\|_2 + \|\widehat{\pi} - \pi\|_{t_n} = o_{P}(1)$ and  $\| \widehat{\mu} - \mu \|_{t_n} =  o_{P}(1)$.
        \item[$(iii)$] $ \| \widehat{\mu} - \mu \|_{t_n} \Big( \|\widehat{\pi} - \pi\|_{t_n} + \|\widehat{\pi} - \pi\|_2 \Big)   = o_{P}\left((n(1-t_n))^{-1/2}\right)$
    \end{enumerate}
    then as $n \to \infty$, it follows that
    \begin{align*}
        \frac{\sqrt{n}}{\sigma_\nu(t_n)}\left(\widehat{\psi}_\nu(t_n) - \psi_\nu(t_n) \right)\overset{d}{\longrightarrow}N(0,1),
    \end{align*}
where $\sigma^2_\nu(t_n) = \mathbb{V}(\varphi_\nu(Z;t_n))$ is the nonparametric efficiency bound defined in \cref{var_corollary}. 
\end{theorem}

\begin{proof}[\textbf{Proof:}] 
     To begin, we derive the kernel identity used in the theorem statement. Then, we split the proof into oracle, empirical process, and remainder sections similar to the proof of \cref{wasserstein_asymptotics}. Again, we assume all nuisance functions are cross-fit, so the rate conditions should be understood to hold foldwise, although we suppress this notation. \medskip

     \textbf{Kernel Identity:} To begin, let $F_t(a) = (1 - t)a + t a_\star$. Then, for $a \in \mathcal{A}_t$, set $u=F_t^{-1}(a)$ and $z=u-a_\star=\frac{a-a_\star}{1-t}$. Next, recall that we define
     \begin{align*}
           \eta_{k,t}(x,a) = \frac{\mathbb{I}(a\in\mathcal{A}_t)}{(1-t)^{k+1}} \sum_{j=0}^{k} \binom{k}{j} \frac{k!}{j!} \left( \frac{a-a_\star}{1-t}\right)^j \pi^{(j)}\left(\frac{a-ta_\star}{1-t}\mid x\right)
     \end{align*}
     such that $\eta_{0, t} = \nu_t$. Then, observe that
     \begin{align*}
     \omega_{p,t}(x,a) &= \sum_{k=0}^{p}\frac{(1-t)^k}{k!}\eta_{k,t}(x,a)\\ 
     &= \frac{\mathbb{I}(a\in\mathcal{A}_t)}{1-t} \sum_{k=0}^{p}\sum_{j=0}^{k} \binom{k}{j}\frac{z^j}{j!}\pi^{(j)}(u\mid x)\\
     &= \frac{\mathbb{I}(a\in\mathcal{A}_t)}{1-t} \sum_{j=0}^{p} \left(\sum_{k=j}^{p}\binom{k}{j}\right) \frac{(u-a_\star)^j}{j!}\pi^{(j)}(u\mid x)\\
     &= \frac{\mathbb{I}(a\in\mathcal{A}_t)}{1-t}K_p(u\mid x),
\end{align*}
where the last equality uses the identity $ \sum_{k=j}^{p}\binom{k}{j} = \binom{p+1}{j+1}$. 
Thus, it follows that
\begin{align} \label{omega_kernel_def}
    \omega_{p,t}(x,a) = \frac{\mathbb{I}(a\in\mathcal{A}_t)}{1-t} K_p\!\left(F_t^{-1}(a)\mid x\right).
\end{align}

\textbf{Central Limit Theorem:} In this section, our goal is to establish asymptotic normality of the oracle term $(P_n - P) \{\Phi_{\nu, p}(Z; t)\}$. First, we establish limiting behavior of the variance. Let $\varepsilon = Y-\mu(X,A)$ and $\sigma^2(x,a)=\mathbb{V}(Y\mid X=x,A=a)$. Since $\mathbb{E}(\varepsilon \mid X,A)=0$, the residual term is uncorrelated with every function of $(X,A)$. Consequently,
\begin{align*}
    \sigma_{\nu,p}^2(t_n) = \mathbb{V}(\Phi_{\nu, p}(Z; t_n)) = \mathbb{E} \! \left[ \left(\frac{\omega_{p,t_n}(X,A)}{\pi(A\mid X)} \right)^2\sigma^2(X,A) \right] + \mathbb{V}\!\left(\Gamma_{p,t_n}(X,A)\right).
\end{align*}
Using \cref{omega_kernel_def} and the change of variables $a=F_n(u)=a_\star+(1 - t_n)(u-a_\star)$,
\begin{align*}
    (1 - t_n) \mathbb{E}\!\left[\left(\frac{\omega_{p,t_n}(X,A)}{\pi(A\mid X)} \right)^2\sigma^2(X,A) \right] = \mathbb{E}\!\left[ \int_{\mathcal{A}} K^2_p(u\mid X) \left(\frac{\sigma^2(X,F_n(u))}{\pi(F_n(u)\mid X)} \right) du \right].
\end{align*}
Recall that we assume for $P_X$-almost every $x$, $a\mapsto\pi(a\mid x)$ and $a\mapsto \mathbb{V}(Y\mid X=x,A=a)$ are continuous at $a_\star$. Furthermore, we assume $\sigma_\star^2(X)=\mathbb{V}(Y\mid X,A=a_\star)\geq\sigma_{\min}^2>0$ almost surely. Thus, following the same arguments as in the proof of \cref{var_corollary}, it follows that, as $t_n \to 1$,
\begin{align*}
    (1 - t_n) \mathbb{E}\!\left[\left(\frac{\omega_{p,t_n}(X,A)}{\pi(A\mid X)} \right)^2\sigma^2(X,A) \right] \longrightarrow \underbrace{\mathbb{E} \! \left[ \frac{\sigma_\star^2(X)}{\pi_\star(X)} \| K_p(\, \cdot \mid X) \|^2_{L_2(\mathcal{A})} \right]}_{ V_{p,\star}}.
\end{align*}
Furthermore, we can easily obtain scaling of $\sigma_{\nu,p}^2(t_n)$ as $t_n \to 1$ by noting that $V_{p,\star} \in (0, \infty)$ and $ \mathbb{V}\!\left(\Gamma_{p,t_n}(X,A)\right) = O(1)$. Thus,
\begin{align*}
    \sigma_{\nu,p}^2(t_n) \sim \frac{V_{p,\star}}{1 - t_n} \asymp (1 - t_n)^{-1}.
\end{align*}
Next, we prove the oracle central limit theorem. By condition $(ii)$, there exist $\delta>0$, $C<\infty$, and a measurable
$M_\mu\in L_{2+\delta}(P_X)$ such that, $P_X$-almost surely,
\begin{align*}
    \max_{0\leq k\leq p+1} \sup_{a\in\mathcal{A}}|\mu^{(k)}(X,a)| \leq M_\mu(X), \qquad \max_{0\leq j\leq p} \sup_{a\in\mathcal{A}}|\pi^{(j)}(a\mid X)| \leq C,
\end{align*}
and
\begin{align*}
    \sup_{a\in\mathcal{A}_{t_0}} \mathbb{E}\!\left[ |Y-\mu(X,A)|^{2+\delta} \mid X,A=a \right] \leq C.
\end{align*}
Now, we verify Lindeberg's condition. For each $n$, let $\Phi_{n,i} = \Phi_{\nu,p}(Z_i;t_n)$ and define
\begin{align*}
    s_n^2 = \sum_{i=1}^{n}\mathbb{V}(\Phi_{n,i}) = n\sigma_{\nu,p}^2(t_n).
\end{align*}
For each fixed $n$, the variables $\Phi_{n,1},\ldots,\Phi_{n,n}$ are independent and identically distributed, have mean zero, and have variance $\sigma_{\nu,p}^2(t_n)$. We will show that, for every $\varsigma > 0$,
\begin{align*}
      \frac{1}{s_n^2} \sum_{i=1}^{n} \mathbb{E}\!\left[\Phi_{n,i}^2 \mathbb{I}\!\left( |\Phi_{n,i}|>\varsigma s_n \right) \right] \longrightarrow 0.
\end{align*}
From here, write $\theta_{p,n} = P \{\Gamma_{p,t_n}\}$ and
\begin{align*}
      H_{p,n}(x,a) = \frac{\omega_{p,t_n}(x,a)}{\pi(a\mid x)}
\end{align*}
so that $\Phi_{n,i} = H_{p,n}(X_i,A_i)\varepsilon_i + \Gamma_{p,t_n}(X_i,A_i) - \theta_{p,n}$, where $\varepsilon_i=Y_i-\mu(X_i,A_i)$. Now, let $q = 2+\delta$ for some $\delta > 0$, as supplied by the moment-envelope condition. We first establish a $q$th-moment bound with the correct dependence on $t_n$. By \cref{omega_kernel_def},
\begin{align*}
     H_{p,n}(x,a) = \frac{\mathbb{I}(a\in\mathcal A_{t_n})}{1-t_n} \frac{K_p\!\left(F_n^{-1}(a)\mid x\right)}{\pi(a\mid x) }.
\end{align*}
Let $m_q(x,a) = \mathbb{E}\!\left[ |\varepsilon|^q \mid X=x,A=a \right]$. Then, using the conditional distribution of $A$ given $X$ and the change of
variables $a=F_n(u)=a_\star+(1-t_n)(u-a_\star)$ and $da=(1-t_n)du$, we obtain
\begin{align*}
    \mathbb{E}\left[ \left| H_{p,n}(X,A)\varepsilon \right|^q \right] &= \frac{1}{(1-t_n)^q} \mathbb{E}\!\left[ \int_{\mathcal{A}_{t_n}} \frac{ \left| K_p(F_n^{-1}(a)\mid X) \right|^q m_q(X,a)}{\pi(a\mid X)^{q-1}} da \right] \\
    &= \frac{1}{(1-t_n)^{q-1}} \mathbb{E}\!\left[
        \int_{\mathcal{A}} \frac{ |K_p(u\mid X)|^q m_q(X,F_n(u))}{ \pi(F_n(u)\mid X)^{q-1} } du \right].
\end{align*}
Under our assumed moment-envelope conditions, the expectation on the right-hand side is uniformly bounded. Hence
\begin{align*}
    \mathbb{E}\left[ \left| H_{p,n}(X,A)\varepsilon \right|^q \right] \lesssim
    (1-t_n)^{-(q-1)}.
\end{align*}
The same envelope argument used above for the second moment of
$\Gamma_{p,t}$ gives
\begin{align*}
    \sup_{t\in[t_0,1)} \mathbb{E} \left[ \left| \Gamma_{p,t}(X,A) \right|^q \right] <\infty.
\end{align*}
Then, by Jensen's inequality, $|\theta_{p,n}|^q = |P\{\Gamma_{p,t_n}\}|^q \leq P \{|\Gamma_{p,t_n}|^q \} = O(1)$. Thus, by the elementary inequality $|x+y+z|^q\lesssim |x|^q+|y|^q+|z|^q$, we have that
\begin{align*}
       \mathbb{E}\left[ |\Phi_{n,i}|^q \right] \lesssim (1-t_n)^{-(q-1)}.
\end{align*}
From here, note that on the event $\{|\Phi_{n,1}|> \varsigma s_n\}$ it follows that
\begin{align*}
    |\Phi_{n,1}|^2 \leq \frac{ |\Phi_{n,1}|^q }{(\varsigma s_n)^{q-2}}.
\end{align*}
Therefore,
\begin{align*}
    \frac{1}{s_n^2} \sum_{i=1}^{n} \mathbb{E}\!\left[ \Phi_{n,i}^2 \mathbb{I}\!\left( |\Phi_{n,i}|>\varsigma s_n \right) \right] \leq  \frac{ \mathbb{E}\left[ |\Phi_{n,1}|^q \right]}{\varsigma^{q-2}s_n^{q-2} \sigma_{\nu,p}^2(t_n)} =  \frac{ \mathbb{E}\left[ |\Phi_{n,1}|^q \right]}{ \varsigma^{q-2} n^{(q-2)/2} \sigma_{\nu,p}^q(t_n)}.
\end{align*}
Thus, putting everything together, we have that
\begin{align*}
    \frac{ \mathbb{E}\left[ |\Phi_{n,1}|^q \right]}{ n^{(q-2)/2} \sigma_{\nu,p}^q(t_n)} \lesssim \frac{ (1-t_n)^{-(q-1)} }{ n^{(q-2)/2} (1-t_n)^{-q/2} } = \left\{ n(1-t_n) \right\}^{-(q-2)/2} \longrightarrow 0,
\end{align*}
which follows under the assumption that $n(1-t_n) \to \infty$. Thus, Lindeberg's condition holds and consequently,
\begin{align*}
     \frac{\sqrt{n}}{\sigma_{\nu,p}(t_n)} (P_n-P) \left\{ \Phi_{\nu,p}(Z;t_n) \right\} \overset{d}{\longrightarrow} N(0,1).
\end{align*}


\medskip 

\textbf{Empirical Process Term:} We next consider $(P_n-P) \{  \widehat{\phi}_{\nu, p}(Z; t_n) -\phi_{\nu, p}(Z; t_n) \}$. Again, we assume all nuisance functions are cross-fitted. Then, conditional on the corresponding training sample, they may be regarded as fixed functions that are independent of the validation observations. Since the number of folds is fixed, it suffices to prove the following bounds within a generic validation fold. Thus, we suppress the fold index notation. Now, we define $\Delta_\mu=\widehat\mu-\mu$, $\Delta_H=\widehat H_{p,n}-H_{p,n}$, and $ \Delta_\Gamma = \widehat\Gamma_{p,t_n}-\Gamma_{p,t_n}.$ After grouping like terms, we have the exact decomposition $\widehat{\phi}_{\nu, p}(Z; t_n) -\phi_{\nu, p}(Z; t_n) = D_{1,n}(Z) + D_{2,n}(Z) + D_{3,n}(Z) + D_{4,n}(Z)$ where
\begin{align*}
    D_{1, n}(Z) &= \Delta_H(X,A)(Y-\mu(X,A)), \\ 
    D_{2, n}(Z) &= - H_{p,n}(X,A)\Delta_\mu(X,A),  \\
    D_{3, n}(Z) &= \Delta_\Gamma(X,A), \\
    D_{4, n}(Z) &= - \Delta_H(X,A)\Delta_\mu(X,A).
\end{align*}
The first three terms are first order in the nuisance-estimation errors and can be controlled through their conditional second moments. The fourth term is a second-order product, which we handle separately through its conditional first moment. We begin with bounding the $L_2$ norm of $D_{1,n}$. Since $\mathbb{E}[Y-\mu(X,A)\mid X,A]=0$, it follows that
\begin{align*}
    \|D_{1,n}\|_2^2 = \mathbb{E}\!\left[ \Delta_H(X,A)^2 \sigma^2(X,A) \right] \lesssim \|\Delta_H\|_2^2.
\end{align*}
Next, because
\begin{align*}
      |H_{p,n}(X,A)| \lesssim \frac{\mathbb{I}(A\in\mathcal A_{t_n})}{1-t_n},
\end{align*}
it follows that
\begin{align*}
    \|D_{2,n}\|_2^2 &= \mathbb{E}\!\left[H_{p,n}(X,A)^2 \Delta_\mu(X,A)^2 \right] \\
    &\lesssim \frac{1}{(1-t_n)^2} \mathbb{E}\!\left[ \Delta_\mu(X,A)^2 \mathbb{I}(A\in\mathcal A_{t_n}) \right] \\
    &= \frac{1}{1-t_n} \|\widehat\mu-\mu\|_{t_n}^2.
\end{align*}
Finally, by definition we have $ \|D_{3,n}\|_2 = \|\widehat\Gamma_{p,t_n}-\Gamma_{p,t_n}\|_2.$ Now, define the linear component
\begin{align*}
     g_n(Z) = D_{1,n}(Z) + D_{2,n}(Z) + D_{3,n}(Z).
\end{align*}
Then, applying the triangle inequality and our previously established bounds, it immediately follows that
\begin{align*}
    \sqrt{1-t_n}\,\|g_n\|_2 \lesssim \sqrt{1-t_n}\, \left\| \frac{\widehat\omega_{p,t_n}}{\widehat\pi} - \frac{\omega_{p,t_n}}{\pi} \right\|_2 + \|\widehat\mu-\mu\|_{t_n} + \sqrt{1-t_n}\, \|\widehat\Gamma_{p,t_n}-\Gamma_{p,t_n}\|_2,
\end{align*}
which we assume is $o_P(1)$. Again, we let $\mathcal{T}_n$ denote the training sample associated with the generic validation fold. Conditional on $\mathcal T_n$, the function $g_n$ is fixed and independent of the validation observations. Therefore,
\begin{align*}
    \mathbb{E}\!\left[\left(\frac{\sqrt{n}}{\sigma_{\nu,p}(t_n)} (P_n-P)g_n \right)^2 \mid \mathcal{T}_n \right] =  \frac{\mathbb{V}(g_n(Z)\mid\mathcal{T}_n)}{\sigma_{\nu,p}^2(t_n)} \leq \frac{\|g_n\|_2^2}{\sigma_{\nu,p}^2(t_n)} \lesssim (1-t_n)\|g_n\|_2^2 = o_P(1),
\end{align*}
where the penultimate inequality uses $\sigma_{\nu,p}^2(t_n)\asymp(1-t_n)^{-1}$. The conditional Chebyshev inequality now yields
\begin{align*}
    \frac{\sqrt{n}}{\sigma_{\nu,p}(t_n)} (P_n-P)g_n = o_P(1).
\end{align*}
From here, all that is left is to control $D_{4,n}$. Applying the Cauchy--Schwarz inequality, we find that
\begin{align*}
    P\{|D_{4,n}|\} &= \mathbb{E}\!\left[ |\Delta_H(X,A)\Delta_\mu(X,A)| \right] \\
    &\leq \mathbb{E}\!\left[ \Delta_\mu(X,A)^2 \mathbb{I}(A\in\mathcal A_{t_n})\right]^{1/2} \|\Delta_H\|_2 \\
    &= \|\widehat\mu-\mu\|_{t_n} \left[ \sqrt{1-t_n} \left\| \frac{\widehat\omega_{p,t_n}}{\widehat\pi} - \frac{\omega_{p,t_n}}{\pi} \right\|_2 \right] \\
    &= o_P\!\left\{ (n(1-t_n))^{-1/2} \right\}.
\end{align*}
Again conditioning on $\mathcal T_n$, note the elementary inequality $\mathbb{E}\!\left[ |(P_n-P)f| \mid  \mathcal{T}_n \right] \leq 2P|f|$ holds for every fixed integrable function $f$. Applying this inequality with $f=D_{4,n}$ gives
\begin{align*}
    \mathbb{E}\!\left[ \left| \frac{\sqrt{n}}{\sigma_{\nu,p}(t_n)} (P_n-P)D_{4,n} \right| \mid \mathcal{T}_n \right] \leq \frac{2\sqrt n}{\sigma_{\nu,p}(t_n)} P\{|D_{4,n}| \} \lesssim \sqrt{n(1-t_n)}\,P\{|D_{4,n}|\} = o_P(1),
\end{align*}
where we again used $\sigma_{\nu,p}^2(t_n)\asymp(1-t_n)^{-1}$.
Conditional Markov's inequality therefore implies
\begin{align*}
     \frac{\sqrt{n}}{\sigma_{\nu,p}(t_n)} (P_n-P)D_{4,n} = o_P(1).
\end{align*}
Putting everything together, it now follows that
\begin{align*}
      \frac{\sqrt{n}}{\sigma_{\nu,p}(t_n)} (P_n-P) \left\{ \widehat{\phi}_{\nu, p}(Z; t_n) -\phi_{\nu, p}(Z; t_n) \right\} = o_P(1).
\end{align*}
\medskip 

\textbf{Remainder Bound:} Finally, we must control the remainder of the von Mises decomposition. Observe that
\begin{align*}
      R_{p,n} = \mathbb{E}\!\left[ \widehat\phi_{\nu,p}(Z;t_n) - \phi_{\nu,p}(Z;t_n)\right] = \mathbb{E}\!\left[ \widehat\phi_{\nu,p}(Z;t_n) - \theta_{p, n}\right]
\end{align*}
because $\mathbb{E}\!\left[\phi_{\nu,p}(Z;t_n) \right] = \mathbb{E}\!\left[ \Gamma_{p,t_n}(X,A) \right] = \theta_{p,n}$. Next, let $ \Delta_\mu(x,a) = \widehat\mu(x,a)-\mu(x,a)$ and note that by iterated expectation 
\begin{align*}
    \mathbb{E}\!\left[\widehat H_{p,n}(X,A)(Y-\widehat\mu(X,A))\right] = - \mathbb{E}\!\left[ \widehat H_{p,n}(X,A) \Delta_\mu(X,A)\right].
\end{align*}
Similarly, because $\mathbb{E}\!\left[H_{p,n}(X,A)(Y-\mu(X,A))\right] = 0$, it follows that
\begin{align*}
    R_{p,n} = - \mathbb{E}\!\left[ \widehat H_{p,n}(X,A) \Delta_\mu(X,A)\right] + \mathbb{E}\!\left[ \widehat\Gamma_{p,t_n}(X,A) - \Gamma_{p,t_n}(X,A) \right].
\end{align*}
Next, recall that 
\begin{align*}
     \Gamma_{p,t_n}(x,a) = \sum_{k=0}^{p} \frac{(1-t_n)^k}{k!} (a_\star-a)^k \mu^{(k)} \!\left(x,(1-t_n)a + t_n a_\star \right).
\end{align*}
Under the assumption that the derivative estimators are obtained from the same fitted function $\widehat\mu$, we may write $\widehat\mu^{(k)}-\mu^{(k)} = (\widehat\mu-\mu)^{(k)} = \Delta_\mu^{(k)}$. Thus,
\begin{align*}
    \mathbb{E}\!\left[ \widehat\Gamma_{p,t_n}(X,A) - \Gamma_{p,t_n}(X,A)\right] = \sum_{k=0}^{p} \frac{(1-t_n)^k}{k!} \mathbb{E}\!\left[(a_\star-A)^k \Delta_\mu^{(k)}\!\left( X,(1-t_n)A+t_na_\star\right) \right].
\end{align*}
From here, to rewrite each term in this sum, we fix $x$ and again let $F_{t_n}(a) = (1-t_n)a+t_na_\star$. For $u\in\mathcal{A}_{t_n}$ we define
\begin{align*}
     Q_{k,t_n}(u\mid x) = \left(\frac{a_\star-u}{1-t_n} \right)^k \nu_{t_n}(u\mid x).
\end{align*}
Applying the change of variables $u=F_{t_n}(a)$ we can then see that 
\begin{align*}
    \int_{\mathcal{A}} (a_\star-a)^k \Delta_\mu^{(k)} \!\left(x,F_{t_n}(a)\right) \pi(a\mid x) da &= \int_{\mathcal{A}_{t_n}} \Delta_\mu^{(k)}(x,u) Q_{k,t_n}(u\mid x)du \\
    &\overset{(i)}{=}  \int_{\mathcal{A}} \Delta_\mu(x,u) \eta_{k,t_n}(x,u)du,
\end{align*}
where $(i)$ follows after applying integration by parts $k$ times, and using the fact that
\begin{align*}
     \eta_{k,t_n}(x,u) = (-1)^k \frac{\partial^k}{\partial u^k} Q_{k,t_n}(u\mid x)
\end{align*}
as shown in the proof of \cref{derivative_EIF}. Consequently, it follows that
\begin{align*}
    \mathbb{E}\!\left[ \widehat\Gamma_{p,t_n}(X,A) - \Gamma_{p,t_n}(X,A)\right] &= \mathbb{E}\!\left[\int_{\mathcal{A}} \Delta_\mu(X,u) \left(\sum_{k=0}^{p} \frac{(1-t_n)^k}{k!} \eta_{k,t_n}(X,u) \right) du \right] \\
    &=  \mathbb{E}\!\left[ \int_{\mathcal{A}} \Delta_\mu(X,u) \omega_{p,t_n}(X,u) du \right] \\
    &=  \mathbb{E}\!\left[\frac{\omega_{p,t_n}(X,A)}{\pi(A\mid X)} \Delta_\mu(X,A) \right] \\
    &= \mathbb{E}\!\left[ H_{p,n}(X,A) \Delta_\mu(X,A) \right].
\end{align*}
Then, substituting this result into our definition of $R_{p, n}$ yields
\begin{align*}
    R_{p,n} &= - \mathbb{E}\!\left[\left(\widehat H_{p,n}(X,A)-H_{p,n}(X,A)\right) \Delta_\mu(X,A) \right] \\
    &= \mathbb{E}\!\left[(\mu(X,A)-\widehat\mu(X,A))\left(\frac{\widehat\omega_{p,t_n}(X,A)}{\widehat\pi(A\mid X)} - \frac{\omega_{p,t_n}(X,A)}{\pi(A\mid X)} \right)\right].
\end{align*}
Now, all that remains is to establish a bound on $R_{p,n}$. Note that by construction, $\omega_{p,t_n}$ and $\widehat\omega_{p,t_n}$ both vanish outside
$\mathcal{A}_{t_n}$. Thus, applying the Cauchy--Schwarz inequality, it follows that
\begin{align*}
     |R_{p,n}| &\leq \mathbb{E}\!\left[ \Delta_\mu(X,A)^2 \mathbb{I}(A\in\mathcal A_{t_n})\right]^{1/2} \mathbb{E}\!\left[\left(\widehat H_{p,n}(X,A)-H_{p,n}(X,A) \right)^2 \right]^{1/2} \\
    &= \|\widehat\mu-\mu\|_{t_n} \left(\sqrt{1-t_n} \left\| \frac{\widehat\omega_{p,t_n}}{\widehat\pi} - \frac{\omega_{p,t_n}}{\pi} \right\|_2 \right).
\end{align*}
Finally, because $\sigma_{\nu,p}^2(t_n)\asymp(1-t_n)^{-1}$,
\begin{align*}
    \frac{\sqrt{n}}{\sigma_{\nu,p}(t_n)} |R_{p,n}| &\lesssim
    \sqrt{n(1-t_n)}|R_{p,n}| =o_P(1),
\end{align*}
thereby completing the proof of the remainder bound. 

\medskip 

\textbf{Sharp Intervention:} Now, we consider the approximation argument for the sharp intervention. Recall that
\begin{align*}
    \theta_{p,n} &=\mathbb{E}\!\left[\Gamma_{p,t_n}(X,A)\right] \\
    &= \sum_{k=0}^{p} \frac{(1-t_n)^k}{k!} \mathbb{E}\!\left[ (a_\star-A)^k \mu^{(k)} \!\left(X,(1-t_n)A+t_na_\star\right) \right] \\
    &= \sum_{k=0}^{p} \frac{(1-t_n)^k}{k!} \psi_\nu^{(k)}(t_n).
\end{align*}
Thus, $\theta_{p,n}$ is the order-$p$ Taylor approximation to $\psi_\nu(1)=\psi_\star$, expanded around $t_n$. Then, observe that under our smoothness envelope assumption,
\begin{align*}
    \left|\psi_\nu^{(p+1)}(s)\right| &= \left| \mathbb{E}\!\left[ (a_\star-A)^{p+1} \mu^{(p+1)}\!\left(X,(1-s)A+sa_\star\right) \right] \right| \leq \underbrace{\sup_{a\in\mathcal A}|a_\star-a|^{p+1}\mathbb{E}[M_\mu(X)]}_{C_{p+1}} <\infty.
\end{align*}
The assumed smoothness of $\mu$ along with dominated convergence also gives
\begin{align*}
    \psi_\nu(s) = \mathbb{E}\!\left[\mu\!\left(X,(1-s)A+sa_\star\right) \right] \longrightarrow \mathbb{E}[\mu(X,a_\star)] = \psi_\star.
\end{align*}
Consequently, Taylor's theorem with an integral remainder yields
\begin{align*}
    \psi_\star-\theta_{p,n} = \frac{1}{p!}
    \int_{t_n}^{1} (1-s)^p\psi_\nu^{(p+1)}(s)ds.
\end{align*}
Then, for all sufficiently large $n$,
\begin{align*}
    |\theta_{p,n}-\psi_\star| &\leq \frac{1}{p!} \int_{t_n}^{1} (1-s)^p \left|\psi_\nu^{(p+1)}(s)\right|ds \leq \frac{C_{p+1}}{p!} \int_{t_n}^{1}(1-s)^p ds = \frac{C_{p+1}}{(p+1)!} (1-t_n)^{p+1}.
\end{align*}
Since $\sigma_{\nu,p}^2(t_n)\asymp(1-t_n)^{-1}$, it follows that
\begin{align*}
    \frac{\sqrt{n}}{\sigma_{\nu,p}(t_n)} |\theta_{p,n}-\psi_\star| \lesssim \sqrt{n(1-t_n)}(1-t_n)^{p+1} = \left( n(1-t_n)^{2p+3} \right)^{1/2} \longrightarrow 0.
\end{align*}
With this approximation-bias result complete, we now combine all parts to conclude the proof.

\medskip
\noindent\textbf{Conclusion:}
We now combine the oracle, empirical-process, remainder, and approximation-bias arguments. To now make the aggregation across cross-fitting folds explicit, let $\mathcal{I}_1,\ldots,\mathcal{I}_J$ denote the validation folds, where $J$ is fixed, and write $n_j=|\mathcal{I}_j|$. Let $\mathcal{T}_{n,j}$ denote the training sample for fold $j$, and let $\widehat\phi_{\nu,p}^{(j)}$ be the score constructed from the nuisance functions fitted on $\mathcal{T}_{n,j}$. Define
\begin{align*}
     d_{n,j}(z) = \widehat\phi_{\nu,p}^{(j)}(z;t_n) - \phi_{\nu,p}(z;t_n),
\end{align*}
and $R_{p,n}^{(j)} = \mathbb{E}\!\left[d_{n,j}(Z) \mid \mathcal{T}_{n,j}\right]$, where $Z$ in the conditional expectation is a fresh observation from the true distribution, independent of the training sample. Thus, $R_{p,n}^{(j)}$ is the fold-specific version of the remainder controlled above. Then, define the full empirical-process error and averaged remainder as
\begin{align*}
    \mathcal{E}_{p,n} = \frac{1}{n} \sum_{j=1}^{J} \sum_{i\in\mathcal{I}_j}\left( d_{n,j}(Z_i)-R_{p,n}^{(j)} \right) \qquad \text{and} \qquad \overline{R}_{p,n} = \sum_{j=1}^{J} \frac{n_j}{n}R_{p,n}^{(j)}.
\end{align*}
Because $\widehat\psi_{\nu,p}^{\,\mathrm{BC}}(t_n) =\frac{1}{n} \sum_{j=1}^{J} \sum_{i\in\mathcal{I}_j} \widehat\phi_{\nu,p}^{(j)}(Z_i;t_n)$ and
$\Phi_{\nu,p}(Z;t_n)=\phi_{\nu,p}(Z;t_n)-\theta_{p,n}$, we have the exact decomposition
\begin{align*}
    \widehat\psi_{\nu,p}^{\,\mathrm{BC}}(t_n)-\psi_\star = \frac{1}{n} \sum_{i=1}^{n}\Phi_{\nu,p}(Z_i;t_n) + \mathcal{E}_{p,n} + \overline{R}_{p,n} + \theta_{p,n}-\psi_\star.
\end{align*}
The preceding empirical-process and remainder bounds imply
\begin{align*}
    \frac{\sqrt{n}}{\sigma_{\nu,p}(t_n)} \mathcal{E}_{p,n} = o_P(1),\qquad \frac{\sqrt{n}}{\sigma_{\nu,p}(t_n)} \overline{R}_{p,n} = o_P(1).
\end{align*}
Therefore, it then follows that
\begin{align*}
    \frac{\sqrt{n}}{\sigma_{\nu,p}(t_n)} \left( \widehat\psi_{\nu,p}^{\,\mathrm{BC}}(t_n)-\psi_\star \right) = \frac{1}{\sqrt{n}\sigma_{\nu,p}(t_n)} \sum_{i=1}^{n}\Phi_{\nu,p}(Z_i;t_n) + o_P(1) \overset{d}{\longrightarrow} N(0,1),
\end{align*}
where the final step follows from the oracle central limit theorem and Slutsky's theorem. Finally, the variance calculation gives
\begin{align*}
     (1-t_n)\sigma_{\nu,p}^2(t_n) \longrightarrow V_{p,\star}\in(0,\infty).
\end{align*}
Therefore,
\begin{align*}
    \sqrt{n(1-t_n)} \left( \widehat\psi_{\nu,p}^{\,\mathrm{BC}}(t_n)-\psi_\star \right) = \sqrt{1-t_n}\sigma_{\nu,p}(t_n) \frac{\sqrt{n}}{\sigma_{\nu,p}(t_n)} \left( \widehat\psi_{\nu,p}^{\,\mathrm{BC}}(t_n)-\psi_\star \right) \overset{d}{\longrightarrow} N(0,V_{p,\star})
\end{align*}
by another application of Slutsky's theorem.

\end{proof}

\end{document}
